% === Кодировки и язык ===
\usepackage{cmap}
\usepackage[T2A]{fontenc}
\usepackage{mathtext}
\usepackage[utf8]{inputenc}
\usepackage[english, russian]{babel}

% === Цвета ===
\usepackage[dvipsnames]{xcolor}

% === Математика ===
\usepackage{amsmath, amssymb, wasysym, amsthm}
\usepackage{stmaryrd}
\usepackage{proof}

% === Графика и float'ы ===
\usepackage{wrapfig, float}
\usepackage{graphicx}

% === TikZ / PGFPlots ===
\usepackage{pgfplots}
\pgfplotsset{compat=newest}
\usepackage{pgfplotstable}
\usepackage{tikz, tikz-3dplot}
\usetikzlibrary{calc, patterns, angles, quotes, intersections, automata, positioning, arrows.meta}

% === Таблицы ===
\usepackage{tabularx}
\usepackage{xltabular}
\renewcommand\tabularxcolumn[1]{m{#1}}

% === Разметка ===
\usepackage{geometry}
\geometry{margin=0.5in}

% === Списки и прочее ===
\usepackage[shortlabels]{enumitem}
\usepackage{verbatim}

% ==========================================
% === ТЕОРЕМНЫЕ ОКРУЖЕНИЯ (amsthm) ===
% ==========================================
\theoremstyle{plain}
\newtheorem{theorem}{Теорема}[section]
\newtheorem{corollary}{Следствие}[section]
\newtheorem{lemma}{Лемма}[section]
\newtheorem{statement}{Утверждение}[section]
\newtheorem{axiom}{Аксиома}[section]
\newtheorem{law}{Закон}[section]

\theoremstyle{definition}
\newtheorem{definition}{Определение}[section]
\newtheorem{example}{Пример}[section]
\newtheorem{problem}{Задача}
\newtheorem{method}{Метод}[section]
\newtheorem{event}{Событие}[section]

\theoremstyle{remark}
\newtheorem{note}{Заметка}[section]
\newtheorem{property}{Свойство}[section]
\newtheorem{properties}{Свойства}[section]
\newtheorem{principle}{Правило}[section]
\newtheorem{person}{Личность}[section]
\newtheorem*{solution}{Решение}

% ==========================================
% === КРАСИВЫЕ РАМКИ (tcolorbox) ===
% ==========================================
\usepackage{tcolorbox}
\tcbuselibrary{skins, breakable, theorems}

% Базовые стили для разных типов блоков
\tcbset{
    % Стиль для теорем, лемм (Синий - строгий)
    thmstyle/.style={
        enhanced, breakable,
        frame hidden, boxrule=0pt,
        colback=blue!5!white, % Фон: 5% синего, 95% белого
        borderline west={2.5pt}{0pt}{blue!70!black}, % Вертикальная линия слева
        before skip=10pt, after skip=10pt,
        left=4pt, right=4pt, top=4pt, bottom=4pt
    },
    % Стиль для доказательств (фиолетовый - процесс рассуждения)
    proofstyle/.style={
        enhanced, breakable,
        frame hidden, boxrule=0pt,
        colback=Plum!5!white,
        borderline west={2.5pt}{0pt}{Plum!80!black},
        before skip=10pt, after skip=10pt,
        left=4pt, right=4pt, top=4pt, bottom=4pt,
        sharp corners
    },
    % Стиль для определений и аксиом (Зеленый - фундаментальный)
    defstyle/.style={
        enhanced, breakable,
        frame hidden, boxrule=0pt,
        colback=ForestGreen!5!white,
        borderline west={2.5pt}{0pt}{ForestGreen!80!black},
        before skip=10pt, after skip=10pt,
        left=4pt, right=4pt, top=4pt, bottom=4pt
    },
    % Стиль для примеров и задач (Оранжевый - практический)
    exstyle/.style={
        enhanced, breakable,
        frame hidden, boxrule=0pt,
        colback=orange!5!white,
        borderline west={2.5pt}{0pt}{orange!90!black},
        before skip=10pt, after skip=10pt,
        left=4pt, right=4pt, top=4pt, bottom=4pt
    },
    % Стиль для заметок и свойств (Серый - нейтральный)
    remstyle/.style={
        enhanced, breakable,
        frame hidden, boxrule=0pt,
        colback=gray!5!white,
        borderline west={2.5pt}{0pt}{gray!70!black},
        before skip=10pt, after skip=10pt,
        left=4pt, right=4pt, top=4pt, bottom=4pt
    }
}

% Применяем стили к существующим окружениям amsthm
% 1. Строгая математика
\tcolorboxenvironment{theorem}{thmstyle}
\tcolorboxenvironment{lemma}{thmstyle}
\tcolorboxenvironment{corollary}{thmstyle}
\tcolorboxenvironment{statement}{thmstyle}
\tcolorboxenvironment{law}{thmstyle}

% 2. Доказательство
\tcolorboxenvironment{proof}{proofstyle}
\renewcommand{\qed}{}

% 3. Определения
\tcolorboxenvironment{definition}{defstyle}
\tcolorboxenvironment{axiom}{defstyle}

% 4. Практика
\tcolorboxenvironment{example}{exstyle}
\tcolorboxenvironment{problem}{exstyle}
\tcolorboxenvironment{method}{exstyle}
\tcolorboxenvironment{event}{exstyle}

% 5. Заметки и прочее
\tcolorboxenvironment{note}{remstyle}
\tcolorboxenvironment{property}{remstyle}
\tcolorboxenvironment{properties}{remstyle}
\tcolorboxenvironment{principle}{remstyle}
\tcolorboxenvironment{person}{remstyle}
\tcolorboxenvironment{solution}{remstyle}

% === Гиперссылки (почти последним!) ===
\usepackage[
colorlinks=true,
linkcolor=blue,
urlcolor=blue,
citecolor=blue,
bookmarks=true,
bookmarksopen=false
]{hyperref}
